\documentclass[11pt,addpoints]{exam}
\usepackage[letterpaper,margin=1in,headheight=14pt]{geometry}
\usepackage{amsmath,amssymb}

\newcommand{\F}{\mathbb{F}}
\newcommand{\R}{\mathbb{R}}
\newcommand{\C}{\mathbb{C}}
\newcommand{\PP}{\mathcal{P}}
\newcommand{\BB}{\mathcal{B}}
\newcommand{\CC}{\mathcal{C}}
\DeclareMathOperator{\Ran}{Ran}
\DeclareMathOperator{\rank}{rank}
\DeclareMathOperator{\tr}{trace}
\DeclareMathOperator{\spn}{span}
\newcommand{\gap}{\vspace{\stretch{1}}}

\pagestyle{headandfoot}
\firstpageheader{}{}{}
\firstpagefooter{}{}{}
\runningheader{MATH 4301/5031, Fall 2026}{}{Supplementary Practice Midterm 1, Page \thepage\ of \numpages}
\runningheadrule
\runningfooter{}{}{}

\pointsinmargin
\bracketedpoints
\pointname{}

\begin{document}

%======================== COVER ========================
\noindent\textbf{\large MATH 4301/5031 \hfill Supplementary Practice Midterm 1 \hfill Fall 2026}

\vspace{1.5cm}
\noindent\textbf{Full Name:} \rule{4.3in}{0.4pt}

\vspace{1.2cm}
\noindent\textbf{\large Instructions}
\begin{enumerate}
  \item \textbf{Part I} (Problems 1--5) follows the format of the practice midterm, with every problem pushed further. Time yourself: 75 minutes. \textbf{Part II} (Problems 6--8 and the Bonus) goes beyond it and is untimed.
  \item Answers should be completely justified unless otherwise noted. Counterexamples must be explicit, and you should check that they work.
  \item Do not use determinants, eigenvalues, or inner products.
  \item Notation: $\F$ denotes $\R$ or $\C$. $\ker T$ and $\Ran T$ are the kernel and range of $T$. $\PP_n$ is the space of real polynomials of degree at most $n$, and $\PP$ is the space of all real polynomials. $E_{ij}$ is the matrix with a $1$ in entry $(i,j)$ and $0$ everywhere else.
  \item Write your answers in the space provided. If you need more space, continue on the back.
\end{enumerate}

\vfill
\begin{center}
\renewcommand{\arraystretch}{1.8}
\begin{tabular}{|l|c|c|c|c|c||c|c|c||c||c|}
\hline
Problem & 1 & 2 & 3 & 4 & 5 & 6 & 7 & 8 & Bonus & Total \\ \hline
Points  & 24 & 18 & 22 & 14 & 22 & 15 & 15 & 20 & $+10$ & 150 \\ \hline
Score   &    &    &    &    &    &    &    &    &       &     \\ \hline
\end{tabular}
\end{center}
\vfill
\begin{center}
\textit{Do not turn the page until you start the clock.}
\end{center}
\newpage

\begin{questions}

%======================== PART I ========================
\item[]\textbf{\large Part I}\par\medskip

%------------------------ Problem 1 ------------------------
\question[24] Determine whether each statement below is true or false. Briefly justify your answer with a proof or an explicit counterexample.
\begin{parts}
\part The map $T\colon \C \to \C$, $T(z) = \overline{z}$, is a linear transformation of $\C$ regarded as a vector space over $\C$.
\gap
\part Row operations preserve the column space of a matrix.
\gap
\part If $v_1, v_2, v_3 \in V$ are linearly dependent, then $v_3$ is a linear combination of $v_1$ and $v_2$.
\gap
\part If $A \in M_{m\times n}(\F)$ and the equation $Ax = b$ has exactly one solution for some $b \in \F^m$, then $Ax = c$ has at most one solution for every $c \in \F^m$.
\gap
\newpage
\part If the columns of $A \in M_{m\times n}(\F)$ are linearly independent, then $A$ is right invertible.
\gap
\part If $X$ and $Y$ are subspaces of $V$ and $X \cup Y = V$, then $X = V$ or $Y = V$.
\gap
\part There exists a linear map $T\colon \R^5 \to \R^5$ with $\ker T = \Ran T$.
\gap
\part Every $4$-dimensional subspace of $\R^6$ and every $3$-dimensional subspace of $\R^6$ have a nonzero vector in common.
\gap
\end{parts}
\newpage

%------------------------ Problem 2 ------------------------
\question[18] Let $V$ be a vector space over $\F$.
\begin{parts}
\part Define what it means for vectors $v_1, \dots, v_m \in V$ to be \emph{linearly independent}, and what it means for them to be a \emph{basis} of $V$.
\gap
\part Using only the axioms of a vector space, prove that $\alpha 0 = 0$ for every $\alpha \in \F$. Then prove that if $\alpha \in \F$ and $v \in V$ satisfy $\alpha v = 0$, then $\alpha = 0$ or $v = 0$.
\gap
\newpage
\part Let $v_1, \dots, v_n$ be a basis of $V$, and for $k = 1, \dots, n$ let $w_k = v_k + v_{k+1} + \dots + v_n$. Prove that $w_1, \dots, w_n$ is a basis of $V$.
\gap
\part Let $n \ge 2$ and let $v_1, \dots, v_n \in V$ be linearly independent. Determine, with proof, all $n$ for which the vectors
\[
v_1 + v_2,\quad v_2 + v_3,\quad \dots,\quad v_{n-1} + v_n,\quad v_n + v_1
\]
are linearly independent.
\gap
\end{parts}
\newpage

%------------------------ Problem 3 ------------------------
\question[22] \mbox{}
\begin{parts}
\part Let $U = \{x \in \F^6 : x_1 + x_2 + x_3 = 0 \text{ and } x_4 = x_5 = x_6\}$. Find the smallest positive integer $m$ for which there exists a linear map $T\colon \F^6 \to \F^m$ with $\ker T = U$. Prove that no smaller $m$ works, and write down such a $T$ explicitly for your value of $m$.
\gap
\part Let $V$ be finite-dimensional and let $S, T\colon V \to V$ be linear. Prove that if $ST$ is invertible, then $TS$ is invertible.
\gap
\newpage
\part Show that the conclusion of (b) can fail when $V$ is infinite-dimensional: find linear maps $S, T\colon \PP \to \PP$ such that $ST = I$ but $TS$ is not invertible.
\gap
\part Now let $V$ be an arbitrary vector space, possibly infinite-dimensional, and let $S, T\colon V \to V$ be linear. Prove that if $STS$ is invertible, then $S$ and $T$ are both invertible.
\gap
\end{parts}
\newpage

%------------------------ Problem 4 ------------------------
\question[14] Let $V$ be the set of all $A \in M_{3\times 3}(\R)$ such that every row and every column of $A$ sums to $0$.
\begin{parts}
\part Show that $V$ is a subspace of $M_{3\times 3}(\R)$.
\gap
\part Find $\dim V$. Justify your answer by exhibiting a basis of $V$ and proving that it is a basis.
\gap
\newpage
\part Let $W$ be the set of all $A \in M_{3\times 3}(\R)$ whose three row sums and three column sums are all equal to one another (the common value may depend on $A$). Prove that $W$ is a subspace of $M_{3\times 3}(\R)$ and that $\dim W = \dim V + 1$.
\gap
\end{parts}
\newpage

%------------------------ Problem 5 ------------------------
\question[22] Let $V = \{X \in M_{2\times 2}(\R) : X^T = X\}$, with basis $\BB = (E_{11},\ E_{12} + E_{21},\ E_{22})$. Let
\[
M = \begin{pmatrix} 1 & 1 \\ 1 & 2 \end{pmatrix},
\]
and define $T\colon V \to V$ by $T(X) = M X M^T$.
\begin{parts}
\part Show that $T$ maps $V$ into $V$ and that $T$ is linear.
\gap
\part Find the matrix $[T]_{\BB}$ of $T$ with respect to $\BB$.
\gap
\part Show that $T$ is invertible, and compute $[T^{-1}]_{\BB}$.
\gap
\newpage
\part Find $X \in V$ such that $T(X) = \begin{pmatrix} 1 & 2 \\ 2 & 3 \end{pmatrix}$.
\gap
\part The matrices
\[
C_1 = \begin{pmatrix} 1 & 0 \\ 0 & 0 \end{pmatrix}, \qquad
C_2 = \begin{pmatrix} 1 & 1 \\ 1 & 0 \end{pmatrix}, \qquad
C_3 = \begin{pmatrix} 1 & 1 \\ 1 & 1 \end{pmatrix}
\]
form a basis $\CC$ of $V$. Write down the change-of-coordinates matrix that converts $\CC$-coordinates into $\BB$-coordinates, and use it to compute $[T]_{\CC}$.
\gap
\end{parts}
\newpage

%======================== PART II ========================
\uplevel{\textbf{\large Part II}\par\medskip}

%------------------------ Problem 6 ------------------------
\question[15] Let $D\colon \PP_3 \to \PP_3$ be differentiation, $Dp = p'$, and let $T = I - D$, so that $(Tp)(t) = p(t) - p'(t)$.
\begin{parts}
\part Write down the matrices of $D$ and of $T$ with respect to the basis $1, t, t^2, t^3$ of $\PP_3$.
\gap
\part Let $N\colon V \to V$ be a linear map on any vector space $V$ such that $N^k = 0$ for some $k \ge 1$. Prove that $I - N$ is invertible, with
\[
(I - N)^{-1} = I + N + N^2 + \dots + N^{k-1}.
\]
Deduce that $T$ is invertible, and write down the matrix of $T^{-1}$ with respect to $1, t, t^2, t^3$.
\gap
\newpage
\part Without row reducing anything, find the unique $p \in \PP_3$ with $p(t) - p'(t) = t^3 - 3t^2 + 1$.
\gap
\part Prove that $p \mapsto p - p'$ is an invertible linear map $\PP \to \PP$, even though differentiation $D\colon \PP \to \PP$ is not invertible.
\gap
\end{parts}
\newpage

%------------------------ Problem 7 ------------------------
\question[15] For a square matrix $X = (x_{ij})$, let $\tr X = \sum_i x_{ii}$.
\begin{parts}
\part Let $A \in M_{m\times n}(\F)$ and $B \in M_{n\times m}(\F)$. Prove that $\tr(AB) = \tr(BA)$. (Note that $AB$ and $BA$ have different sizes when $m \ne n$.)
\gap
\part Prove that there are no matrices $A, B \in M_{n\times n}(\F)$ with $AB - BA = I$.
\gap
\newpage
\part Let $W = \spn\{AB - BA : A, B \in M_{n\times n}(\F)\}$. Prove that
\[
W = \{X \in M_{n\times n}(\F) : \tr X = 0\},
\]
and find $\dim W$.
\gap
\end{parts}
\newpage

%------------------------ Problem 8 ------------------------
\question[20] \mbox{}
\begin{parts}
\part Let $A \in M_{m\times n}(\F)$. Prove that if $Ax = 0$ has only the trivial solution, then $A^T y = c$ has a solution for every $c \in \F^n$.
\gap
\part Let $A \in M_{m\times n}(\F)$ and $B \in M_{n\times p}(\F)$. Prove that
\[
\rank(AB) \ge \rank A + \rank B - n.
\]
(Hint: apply the rank theorem to $A$ restricted to $\Ran B$.)
\gap
\newpage
\part Let $A \in M_{n\times n}(\F)$ satisfy $A^2 = A$. Prove that $\rank A + \rank(I - A) = n$.
\gap
\end{parts}

\end{questions}
\newpage

%------------------------ Bonus ------------------------
\noindent\hspace*{-2.2em}\textbf{[+10]}\hspace{0.9em}\textbf{Bonus.} Prove that a vector space over $\R$ cannot be written as the union of three proper subspaces. (A subspace $U \subseteq V$ is \emph{proper} if $U \ne V$.)
\gap

\end{document}
